% !PRINTED 302
\weberpage{302}

2. $\mu$ odd:

\[
\begin{array}{ll}
\text{a)}\ f^{(\nu-1)}(x),\ f^{(\nu)}(x), & \text{b)}\ f^{(\nu-1)}(x),\ f^{(\nu)}(x) \\[2pt]
\delta_1,\quad + \quad - & \quad - \quad - \\[2pt]
\delta_2,\quad + \quad + & \quad - \quad +
\end{array}
\]

Thus for even $\mu$ no change of sign is lost between $f^{(\nu-1)}$ and $f^{(\nu)}$; for odd $\mu$ a change of sign is lost or one is gained. Hence the loss of changes of sign on passing through $\xi$ in the sequence

\weq{5}{f^{(\nu-1)},\ f^{(\nu)},\ f^{(\nu+1)}\ .\ .\ .\ f^{(\nu+\mu)}}

is equal to $\mu$ for even $\mu$, and to $\mu \pm 1$ for odd $\mu$; it is therefore always an \gesperrt{even number} and never negative.

We can give these results a clear analytic expression, which lets their meaning be recognised better.

As before, we denote by $V(x)$ the number of changes of sign in the sequence (1) for a definite value of $x$, but with the further stipulation that, if for a value of $x$ some of the terms of the sequences vanish, these are simply passed over in counting the changes of sign. We can then regard $V(x)$ as a function of $x$, which however can change only by integers, and is therefore discontinuous for those values of $x$ for which some terms of the sequence (1) vanish. We shall then denote by $V(x-0)$ and $V(x+0)$ the values of the function $V(x)$ immediately before and immediately after such a place. Then the consideration of (4) shows that in all cases

\weq{6}{V(\xi + 0) = V(\xi)}

and that, if $\xi$ is a $\mu$-fold root of $f(x)$,

\weq{7}{V(\xi - 0) = V(\xi) + \mu + 2h,}

where $h$ is a non-negative integer.

From this it follows, then, if for the present we hold to the assumption that for $x = \alpha$ and $x = \beta$ none of the functions (1) vanishes, that $V(\alpha) - V(\beta)$ differs from the sum of all the numbers $\mu$, that is, from the number of the roots lying in the interval $(\alpha, \beta)$, each counted according to the degree of its multiplicity, by an even, non-negative number; and with this Theorem I is thus freed from its restriction.
